\section{INTRODUCTION}

The Amazon basin holds the largest surface freshwater reserve on the planet and plays a key role in
climate regulation and global biodiversity. Monitoring water quality in this region is increasingly
challenging due to vast distances, dense forest cover, complex hydrography, and limited communication
infrastructure. Traditional manual sampling and laboratory analysis, while accurate, are expensive,
logistically difficult, and unable to provide the temporal resolution needed to detect sudden
contamination events (e.g., oil spills or illegal discharges) or track hydrological changes driven
by extreme climate events.

This context of environmental monitoring under energy and connectivity constraints is also
highlighted in the position paper by \citeonline{aral2024}, which outlines directions for neuromorphic
Edge AI in rural scenarios. In response, the convergence of the Internet of Things (IoT) and
Artificial Intelligence (AI) has been proposed as a viable path for large-scale remote monitoring
\cite{rahu2023}. Connected sensors can stream data in real time, enabling more proactive water
resource management. However, continuous transmission of raw data to the cloud faces severe
connectivity limitations in the Amazon, resulting in high energy consumption and unacceptable
latency. A practical alternative is to move intelligence to the edge, where data are processed
locally on the monitoring device.

The critical challenge, as noted by \citeonline{rahu2023}, is that high-performance machine learning
models such as Long Short-Term Memory (LSTM) networks demand computational and energy resources that
often exceed the capacity of low-cost microcontrollers used in battery- or solar-powered sensor
nodes. Although LSTMs perform well in complex hydrological time series modeling \cite{pyo2023}, their
computational cost can reduce the autonomy of remote stations. In this setting, neuromorphic
computing—and specifically SNNs—offers a promising alternative. SNNs process information asynchronously
through sparse spikes, activating only when significant input variations occur. This property can
provide orders-of-magnitude energy savings compared to conventional artificial neural networks (ANNs)
that operate continuously.

This paper evaluates the technical and practical feasibility of SNN-based edge inference for
predicting water-quality parameters. We focus on multimodal prediction of pH, dissolved oxygen, and
turbidity, and on validating the full pipeline in real conditions, from data processing to offline
deployment. The remainder of this paper is organized as follows: Section~2 presents related work and
positions the study; Section~3 details materials and methods, including data, model, training, and
mobile deployment; Section~4 presents results and performance analyses; Section~5 discusses
limitations and implications; and Section~6 concludes the paper and outlines future directions.

%----------------------------------------------------------------------------------------
%   TRABALHOS RELACIONADOS
%----------------------------------------------------------------------------------------
